\section{Gota el\'iptica: curvatura vari\'avel}
\label{sec:elipse}

\subsection{O teste que faltava}

A Se\c{c}\~ao~\ref{sec:balanco} fechou com uma ressalva expl\'icita: o balan\c{c}o
\'e \emph{exato} apenas quando $\kappa$ \'e constante, pois a\'i
$\sigma\kappa\nabla H = \nabla(\sigma\kappa H)$. Com curvatura vari\'avel sobra
um res\'iduo, e a gota circular n\~ao o mede. Esta se\c{c}\~ao mede.

O caso: elipse de semieixos $a=0{,}2357$ e $b=0{,}1179$ --- raz\~ao $2{:}1$ e
\'area \emph{id\^entica} \`a do c\'irculo $R=1/6$ das se\c{c}\~oes anteriores.
A gota n\~ao est\'a em equil\'ibrio; a tens\~ao superficial a relaxa para o
c\'irculo de mesma \'area, de modo que o raio de equil\'ibrio \'e
$R_{\text{eq}}=\sqrt{ab}=1/6$ e o salto final esperado volta a ser
$\Delta p = 6$. Mesma malha, mesmas propriedades, 1200 passos ($t=1{,}2$, cerca
de quatro tempos de relaxa\c{c}\~ao).

\'E tamb\'em o primeiro caso em que a frente \textbf{deforma}, e portanto o
primeiro em que a \textbf{cirurgia dispara} com o solver no circuito: o n\'umero
de marcadores cai de 108 para 102--104 conforme a elipse se contrai no eixo
maior.

\subsection{O defeito que o c\'irculo escondia}

Na primeira corrida, a formula\c{c}\~ao balanceada \textbf{falhou}. A
circularidade sobe normalmente at\'e o passo~900, chega a $0{,}9887$ --- e ent\~ao
\emph{reverte}, caindo a $0{,}9617$, com a velocidade crescendo. O salto de
press\~ao final sai $10{,}99$ contra os $6$ esperados: \SI{83}{\percent} de erro.

\begin{center}
\begin{tikzpicture}
\begin{axis}[
  width=12cm, height=6.4cm,
  xlabel={passo}, ylabel={circularidade $4\pi A/P^{2}$},
  legend pos=south east, grid=major, legend cell align=left,
  ymin=0.83, ymax=1.01,
]
\addplot table[x=passo, y=espalhado]  {\dat{elipse-circularidade.dat}};
\addlegendentry{espalhado}
\addplot table[x=passo, y=balsuave]   {\dat{elipse-circularidade.dat}};
\addlegendentry{balanceado, $\kappa$ suavizado}
\addplot table[x=passo, y=balvizinho] {\dat{elipse-circularidade.dat}};
\addlegendentry{balanceado, $\kappa$ do vizinho}
\end{axis}
\end{tikzpicture}
\end{center}

O instante da falha \'e a pista: ela aparece justamente quando a gota fica quase
circular, isto \'e, quando a for\c{c}a f\'isica que a movia j\'a decaiu. Enquanto
havia um sinal grande, o artefato ficava submerso; ao ficar sozinho, dominou.

A causa estava numa simplifica\c{c}\~ao que a pr\'opria
Se\c{c}\~ao~\ref{sec:discussao} havia registrado como limita\c{c}\~ao: a
curvatura na faceta era a do \textbf{marcador mais pr\'oximo}. No c\'irculo isso
\'e exato, porque $\kappa$ \'e constante. Na elipse \'e um campo
\emph{descont\'inuo}: salta quando a atribui\c{c}\~ao troca de marcador. As
descontinuidades produzem saltos de for\c{c}a que movem a frente, o que muda a
atribui\c{c}\~ao, o que muda a for\c{c}a --- realimenta\c{c}\~ao.

\subsection{Conserto dirigido, e a confirma\c{c}\~ao do diagn\'ostico}

Trocou-se apenas isso: a curvatura na faceta passou a ser a \textbf{m\'edia
ponderada} das curvaturas dos marcadores vizinhos, com o n\'ucleo de Roma de
largura $h$. Nada mais mudou.

Com $\kappa$ suavizado a instabilidade \textbf{desaparece}, e a trajet\'oria
passa a coincidir com a da formula\c{c}\~ao espalhada at\'e a quarta casa:
$0{,}922213$ contra $0{,}922220$ no passo~200; $0{,}998477$ contra $0{,}998491$
no passo~1150. Duas formula\c{c}\~oes de for\c{c}a diferentes produzindo a mesma
din\^amica \'e evid\^encia forte de que ambas est\~ao certas.

\paragraph{Sem regress\~ao no c\'irculo.} A curvatura \'e constante ali, e a
m\'edia ponderada de um valor constante devolve esse valor exatamente: o
resultado da Se\c{c}\~ao~\ref{sec:balanco} se mant\'em --- $\Delta p = 6{,}000001$
e $\max\lvert u\rvert < \num{5e-11}$.

\subsection{Resultados, e o que eles corrigem}

\pgfplotstabletypeset[
  string type,
  columns/grandeza/.style={column name={grandeza}},
  columns/espalhado/.style={column name={espalhado}},
  columns/balvizinho/.style={column name={bal., $\kappa$ vizinho}},
  columns/balsuave/.style={column name={bal., $\kappa$ suave}},
  columns/vof/.style={column name={VOF}},
  every head row/.style={before row=\toprule, after row=\midrule},
  every last row/.style={after row=\bottomrule},
]{\dat{elipse-resultados.dat}}

Tr\^es leituras, e duas delas corrigem impress\~oes deixadas pelas se\c{c}\~oes
anteriores.

\paragraph{1. Na elipse, o balan\c{c}o n\~ao traz vantagem.} Espalhado
\SI{1.27}{\percent}, balanceado \SI{1.52}{\percent}, VOF \SI{1.19}{\percent} ---
os tr\^es dentro de \SI{0.3}{\percent} um do outro. O ganho espetacular da
Se\c{c}\~ao~\ref{sec:balanco} era \emph{espec\'ifico do caso est\'atico com
$\kappa$ constante}, onde o balan\c{c}o exato \'e alcan\c{c}\'avel. Sob
movimento real e curvatura vari\'avel, as duas formula\c{c}\~oes s\~ao
equivalentes dentro da medida.

\paragraph{2. A vantagem estrutural do VOF em massa aparece agora.} A
Se\c{c}\~ao~\ref{sec:balanco} advertiu que a conserva\c{c}\~ao perfeita do
front-tracking balanceado vinha de o escoamento \emph{ser} est\'atico, e que ``em
escoamento de verdade a deriva volta''. Voltou: \num{1.1e-4} e \num{2.3e-4}
contra \num{2.5e-13} do VOF. Este \'e o primeiro teste com movimento real, e ele
confirma a advert\^encia.

\paragraph{3. Os tr\^es m\'etodos concordam na f\'isica.} Os saltos finais ficam
em $5{,}909$, $5{,}924$ e $5{,}929$ para um valor exato de $6$. O d\'eficit comum
de \SI{1.2}{\percent} n\~ao \'e de m\'etodo: \'e de relaxa\c{c}\~ao incompleta ---
em $t=1{,}2$ a circularidade ainda \'e $0{,}9985$, e resta velocidade da ordem de
\num{9e-3}, que contribui press\~ao din\^amica. Que tr\^es discretiza\c{c}\~oes
independentes errem \emph{junto}, e pelo mesmo tanto, \'e o que se espera de um
erro f\'isico compartilhado, e n\~ao de defeito em qualquer uma delas.

\subsection{A li\c{c}\~ao de m\'etodo}

O resultado da gota circular era verdadeiro e continua v\'alido; o que faltava
era saber \emph{sobre o que} ele falava. Um caso em que uma quantidade \'e
constante n\~ao distingue uma implementa\c{c}\~ao que trata bem essa quantidade
de outra que a trata de qualquer jeito --- e aqui a diferen\c{c}a era entre
funcionar e ir \`a instabilidade.

A limita\c{c}\~ao estava registrada por escrito antes de ser testada, o que
tornou o diagn\'ostico r\'apido: bastou olhar para o \'unico ingrediente que
mudava de comportamento entre c\'irculo e elipse. Registrar uma limita\c{c}\~ao
conhecida no momento em que ela \'e aceita paga-se quando ela cobra.
